\documentclass[10pt]{article}

\usepackage[utf8]{inputenc}

\usepackage[T1]{fontenc}

\usepackage{amsmath}

\usepackage{amsfonts}

\usepackage{amssymb}

\usepackage[version=4]{mhchem}

\usepackage{stmaryrd}



\begin{document}

верительных интервалов для параметров $\beta_{i}$ и для проверки гипотез о значениях, к-рые принимает величина $\beta_{i}$. К роме того, появляется возможность найти доверительные интервалы для $\mathrm{E}\left(Y \mid x_{1}, \ldots, x_{k}\right)$ при фиксированных значениях всех регрессионных переменных и доверительные интервалы, содержащие следующее $(n+1)$-е значение величины $y$ (т. н. интервалы предсказания). Наконец, можно на основе вектора выборочных коэффициентов регрессии $\widehat{\beta}$ построить доверительный эллипсоид для вектора $\beta$ или для любой совокупности неизвестных коэффициентов регрессии, а также доверительную область для всей линии или прямой регрессии.



Р. а. является одним из наиболее распространенных методов обработки экспериментальных данвых при изучении зависимостей в физике, биологии, экономике, технике и др. областях. На моделях Р. а. основаны такие разделы математич. статистики, как дисперсионный анализ и планирование әксперимента, эти модели широко используются в многомерном статистическом анализе.



Лuт.: [1] К е н д а л л М. Д ж., С ть ю а р т А., Статистические выводы и связи, пер. с англ., М., 1973; [2] С м и рн о в Н. В., Д у н и н-Б а р к о в с к и й Н. В., Курс теории вероятностей и математической статистики для технических приложений, 3 изд., М., 1969; [3] А й в а з я н С. А., Статистическое исследование зависимостей, М., 1968; [4] Р а о С. Р., Линейные статистические методы и их применения, пер. с англ., М., 1968; [5] Д р е й п е р Н., С м и т Г., Прикладной регрес-

сионный анализ, пер. с англ., М., 1973.

А. В. Прохоров.



РЕГРЕССИЯ - зависимость среднего значения какой-либо случайной величины от нек-рой другой величины или от нескольких величин. Если, например, при каждом значении $x=x_{i}$ наблюдается $n_{i}$ значений $y_{i 1}, \ldots, y_{i n_{i}}$ случайной величины $Y$, то зависимость средних арифметических



\[

\bar{y}_{i}=\frac{1}{n_{i}}\left(y_{i \mathbf{1}}+\ldots+y_{i n_{i}}\right)

\]



этих значений от $x_{i}$ и является Р. в статистич. понимании этого термина. При обнаруженной закономерности изменения $\bar{y}$ с изменением $x$ предполагается, пто в основе наблюдаемого явления лежит вероятностная зависимость: при каждом фиксированном значении $x$ случайная величина $Y$ имеет определенное распределение вероятностей с математич. ожиданием, к-рое является функцией $x$ :



\[

\mathrm{E}(Y \mid x)=m(x)

\]



Зависимость $y=m(x)$, где $x$ играет роль «независимой» переменной, наз. р е г p е с с и е й (или ф у ш к циe й $\mathrm{p}$ е $\mathrm{r} \mathrm{p}$ е c c и и) в вероятностном понимании этого термина. График функции $m(x)$ наз. л и н и е й $\mathrm{p} \mathrm{e}-$ г р е с с и и, или к р и в й ре г ре с с и и, величины $Y$ по $x$. Переменная $x$ наз. $\mathrm{p}$ е г $\mathrm{p}$ е с $\mathrm{c}$ и о н н о й п е р е менно й, или р е г р е с с о ром. Точность, с к-рой линия регрессии $Y$ по $x$ передает изменение $Y$ в среднем при изменении $x$, измеряется дисперсией величины $Y$, вычисляемой для каждого значения $x$ : $\mathrm{D}(Y \mid x)=\sigma^{2}(x)$.



Графически зависимость дисперсии $\sigma^{2}(x)$ от $x$ выражается т. н. ск е д а с тиче с к о й ли ни е й. Если $\sigma^{2}(x)=0$ при всех значениях $x$, то с вероятностью 1 величины связаны строгой функциональной зависимостью. Если $\sigma^{2}(x) \neq 0$ ни при каком значении $x$ и $m(x)$ не зависит от $x$, то регрессия $Y$ по $x$ отсутствует.



В теории вероятностей задача Р. решается применительно к такой ситуации, когда значения регрессионной переменной $x$ соответствуют значениям нек-рой случайной величины $X$ и предполагается известным совместное распределение вероятностей величин $X$ и $Y$ (при этом математич. ожидание $\mathrm{E}(Y \mid x)$ и дисперсия $\mathrm{D}(Y \mid x)$ будут соответственно условным математич. ожиданием и условной дисперсией случайной величины $Y$ при фиксированном значении $X=x)$. В этом случае определены две Р.: $Y$ по $x$ и $X$ по $y$, и понятие Р. может быть использовано также для того, чтобы ввести некрые меры взаимосвязанности случайных величин $X$ и $Y$, определяемые как характеристики степени концентрации распределения около линий Р. (см. Корреялиия).



Функции Р. обладают тем свойством, что среди всех действительных функций $f(x)$ минимум математич. ожидания $\mathrm{E}(Y-f(x))^{2}$ достигается для функции $f(x)=$ $=m(x)$, то есть регрессия $Y$ по $x$ дает наилучшее (в указанном смысле) представление величины $Y$. Наиболее важным является тот случай, когда регрессия $Y$ по $x$ ли н е йн а, т. е.



\[

\mathrm{E}(Y \mid x)=\beta_{0}+\beta_{1}(x) .

\]



Коәффициенты $\beta_{0}$ и $\beta_{1}$, наз. к о э фф и ц и е н т а м и Р., легко вычисляются:



\[

\beta_{0}=m_{Y}-\rho \frac{\sigma_{Y}}{\sigma_{X}} m_{X}, \beta_{1}=\rho \frac{\sigma_{Y}}{\sigma_{X}}

\]



(здесь $\rho$-корреляции коәфбициент $X$ и $Y, m_{X}=\mathrm{E} X$, $\left.m_{Y}=\mathrm{E} Y, \sigma_{X}^{2}=\mathrm{D} X, \sigma^{2}=\mathrm{D} Y\right)$, и п р я м а я р е г р е сс и и $Y$ по $x$ имеет вид



\[

y=m_{Y}+\rho \frac{\sigma_{Y}}{\sigma_{\boldsymbol{X}}}\left(x-m_{\boldsymbol{X}}\right)

\]



(аналогичным образом находится прямая регрессии $X$ по $y$ ). Точная линейная $\mathrm{P}$. имеет место в случае, когда двумерное распределение величин $X$ и $Y$ является нормальным.



В условиях статистич. приложений, когда для точного определения Р. нет достаточных сведений о форме совместного распределения вероятностей, возникает задача приближенного нахождения Р. Решению этой задачи может служить выбор из всех функций $g(x)$, принадлежащих заданному классу, такой функции, к-рая дает наилучшее представление величины $Y$ в том смысле, что минимизирует математич. ожидание $\mathrm{E}(Y-g(X))^{2}$. Найденная функция наз. с р е д н е й к в а д р а т ичес к о й $\mathrm{P}$.



Простейшим будет случай л и н е й н о й с ре д н е й к в а д р а т и ч е с к о й Р., когда отыскивают наилучшую линейную аппроксимацию величины $Y$ посредством величины $X$, т. е. такую линейную функцию $g(x)=\beta_{0}+\beta_{1} x$, для к-рой выражение $\mathrm{E}(Y-g(X))^{2}$ принимает наименьшее возможное значение. Данная экстремальная задача имеет единственное решение



\[

\beta_{0}=m_{Y}-\beta_{1} m_{X}, \quad \beta_{1}=\rho \frac{\sigma_{Y}}{\sigma_{X}},

\]



т. е. вычисление приближенной линии Р. приводит к тому же результату, к-рый получен в случае точной линейной Р.:



\[

y=m_{Y}+\rho \frac{\sigma_{Y}}{\sigma_{X}}\left(x-m_{X}\right)

\]



Минимальное значение $\mathrm{E}(Y-g(X))^{2}$ при вычисленных значениях параметров равно $\sigma_{Y}^{2}\left(1-\rho^{2}\right)$. Если регрессия $m(x)$ существует, то при любых $\beta_{0}$ и $\beta_{1}$ имеет место соотношение



$\mathrm{E}\left[Y-\beta_{0}-\beta_{1} X\right]^{2}=\mathrm{E}[Y-m(X)]^{2}+\mathrm{E}\left[m(X)-\beta_{0}-\beta_{1} X\right]^{2}$, откуда следует, что прямая средней квадратич. регрессии $y=\beta_{0}+\beta_{1} x$ дает наилучшее приближение к линии регрессии $m(x)$, если измерять расстояние вдоль оси $y$. Поэтому если линия $m(x)$ есть прямая, то она совпадает с прямой средней квадратической $\mathrm{P}$.



В общем случае, когда Р. сильно отличается от линейной, можно поставить задачу нахождения многочлена $g(x)=\beta_{0}+\beta_{1} x+\ldots+\beta_{m} x^{m}$ нек-рой стешени $m$, для





\end{document}